\begin{ouverture}[calcul]{9}{Comparer et arrondir les décimaux}{Partie 2 · Fractions et nombres décimaux}
{Ordre · Encadrements · Partie entière · Arrondis}
\objectif{Comparer, encadrer et intercaler des nombres décimaux.}
\objectif{Placer des décimaux sur une demi-droite graduée.}
\objectif{Donner la partie entière et l’arrondi à l’entier.}
\prerequis{Écriture à virgule, dixièmes et centièmes.}
\end{ouverture}
\pagestyle{otMain}\essentiel\fontsize{11}{15}\selectfont
\phantomsection\label{cours-9-1}
\essSection{1}{Comparer et ranger}
\begin{propriete}Pour comparer deux décimaux : comparer d’abord les parties entières ; si elles sont égales, comparer les chiffres des dixièmes, puis ceux des centièmes. Ainsi $6{,}8>6{,}08$, car 8 dixièmes dépassent 0 dixième.\end{propriete}
\begin{exemple}Comparer $3{,}5$ et $3{,}45$ : écrire $3{,}5=3{,}50$. Comme $50$ centièmes dépassent $45$ centièmes, $3{,}5>3{,}45$. Un nombre n’est pas plus grand parce qu’il a plus de chiffres.\end{exemple}
\label{fig-cours-9-1-cours-droite-decimaux}\figureGeo{cours-droite-decimaux}

\begin{savoirfaire}[Intercaler un nombre]Entre deux décimaux, on peut toujours en intercaler un autre : entre $2{,}4$ et $2{,}5$, penser en centièmes : $2{,}40<2{,}45<2{,}50$. Ajouter un rang suffit pour affiner.\end{savoirfaire}
\clearpage
\phantomsection\label{cours-9-2}
\essSection{2}{Partie entière et arrondi}
\begin{definition}La partie entière de $17{,}83$ est 17 : c’est le nombre entier juste avant. On a l’encadrement $17<17{,}83<18$ entre deux entiers consécutifs.\end{definition}
\begin{propriete}L’arrondi à l’entier est l’entier le plus proche. Si le chiffre des dixièmes est 0, 1, 2, 3 ou 4, on garde la partie entière ; sinon on prend l’entier suivant. L’arrondi de $7{,}3$ est 7 ; celui de $7{,}83$ est 8. Pour $12{,}5$, à égale distance, on choisit 13.\end{propriete}
\begin{exemple}Sur la demi-droite, $7{,}83$ est plus près de 8 que de 7 : l’arrondi se voit sur les positions. Un prix de $7{,}83$ euros est proche de 8 euros : l’arrondi aide à estimer.\end{exemple}
\begin{attention}Arrondir n’est pas tronquer : l’arrondi de $9{,}8$ est 10, pas 9. Et deux nombres arrondis au même entier restent différents : l’arrondi sert à estimer, pas à remplacer la valeur exacte.\end{attention}

\clearpage\rubrique{Découverte}
\begin{blocLivre}{ex-9-1}\begin{exercice}[1]{Le piège classique}
Comparer $3{,}5$ et $3{,}45$.
\end{exercice}\end{blocLivre}
\begin{blocLivre}{ex-9-2}\begin{exercice}[1]{Compléter par un symbole}
Compléter par $<$ ou $>$ : $6{,}08\ \ldots\ 6{,}8$, puis $12{,}34\ \ldots\ 9{,}99$.
\end{exercice}\end{blocLivre}
\begin{blocLivre}{ex-9-3}\begin{exercice}[1]{Ranger cinq nombres}
Ranger dans l’ordre croissant : $7{,}3$ ; $7{,}09$ ; $7{,}31$ ; $7{,}1$ ; $6{,}99$.
\end{exercice}\end{blocLivre}
\begin{blocLivre}{ex-9-4}\begin{exercice}[1]{Encadrer par des entiers}
Encadrer chaque nombre entre deux entiers consécutifs : $4{,}62$, puis $0{,}85$.
\end{exercice}\end{blocLivre}
\begin{blocLivre}{ex-9-5}\begin{exercice}[1]{Partie entière}
Donner la partie entière de $17{,}83$, de $9{,}04$ et de $0{,}76$.
\end{exercice}\end{blocLivre}

\rubrique{Entraînement}
\begin{blocLivre}{ex-9-6}\begin{exercice}[2]{Lire entre les graduations}
La demi-droite est graduée de $0{,}1$ en $0{,}1$. Lire les abscisses de A et B, puis placer $C(5{,}45)$.
\label{fig-ex-9-6}\figureGeo{dec9-6}
\end{exercice}\end{blocLivre}
\begin{blocLivre}{ex-9-7}\begin{exercice}[2]{Intercaler}
Donner un nombre compris entre $2{,}4$ et $2{,}5$, puis un nombre compris entre $8{,}31$ et $8{,}32$.
\end{exercice}\end{blocLivre}
\begin{blocLivre}{ex-9-8}\begin{exercice}[2]{Arrondir à l’entier}
Donner l’arrondi à l’entier de $7{,}3$, de $7{,}83$ et de $12{,}5$.
\end{exercice}\end{blocLivre}
\begin{blocLivre}{ex-9-9}\begin{exercice}[2]{Compter de 0,1 en 0,1}
Compter de $0{,}1$ en $0{,}1$ à partir de $3{,}7$, jusqu’à dépasser 4. Écrire les nombres obtenus.
\end{exercice}\end{blocLivre}
\begin{blocLivre}{ex-9-10}\begin{exercice}[2]{Le même nombre ?}
Dire si les écritures désignent le même nombre : $5{,}60$ et $5{,}6$ ; $5{,}06$ et $5{,}6$ ; $07{,}2$ et $7{,}2$.
\end{exercice}\end{blocLivre}

\rubrique{Problèmes}
\begin{blocLivre}{ex-9-11}\begin{exercice}[3]{Le podium du saut}
Aux essais, Lina saute $3{,}45$ m, Maya $3{,}5$ m et Inès $3{,}09$ m. Donner le classement de la meilleure à la moins bonne performance.
\end{exercice}\end{blocLivre}
\begin{blocLivre}{ex-9-12}\begin{exercice}[3]{La toise de la classe}
Quatre élèves mesurent $1{,}38$ m, $1{,}4$ m, $1{,}29$ m et $1{,}33$ m. Les ranger du plus petit au plus grand.
\end{exercice}\end{blocLivre}
\begin{blocLivre}{ex-9-13}\begin{exercice}[3]{Toujours un intermédiaire ?}
Donner deux nombres différents compris entre $6{,}2$ et $6{,}3$. Peut-on toujours trouver un nombre entre deux décimaux différents ? Expliquer.
\end{exercice}\end{blocLivre}
\begin{blocLivre}{ex-9-14}\begin{exercice}[3]{L’étiquette effacée}
Dans un magasin, un paquet coûte $4{,}05$ euros et un autre $4{,}5$ euros. Tom pense que $4{,}05$ est le plus cher « car 05 vient après 5 ». Corriger son raisonnement.
\end{exercice}\end{blocLivre}
\begin{blocLivre}{ex-9-15}\begin{exercice}[3]{L’erreur de rangement}
Aya écrit : « $12{,}34>12{,}5$ car $34>5$. » Expliquer son erreur, donner la bonne comparaison, puis l’arrondi à l’entier des deux nombres.
\end{exercice}\end{blocLivre}

\clearpage
